\documentclass[10pt,a4paper]{article}

\usepackage[T1]{fontenc}
\usepackage[margin=2.5cm]{geometry}
\usepackage{amsmath,amssymb}
\usepackage{booktabs}
\usepackage{array}
\usepackage{float}
\usepackage[per-mode=symbol]{siunitx}
\sisetup{group-minimum-digits=4, group-separator={,}, range-phrase={--}, range-units=single}
\usepackage[hidelinks]{hyperref}
\usepackage{microtype}

% Notation, matching the paper.
\newcommand{\fALL}{\texttt{ALL}}
\newcommand{\fALLgt}{\texttt{ALL>100B}}
\newcommand{\fUL}{\texttt{UL}}
\newcommand{\fULgt}{\texttt{UL>100B}}
\newcommand{\rpi}{Raspberry~Pi}
\newcommand{\jn}{Jetson}

\title{Supplementary Material\\[2mm]
\large Energy-Aware Assessment and Energy Modeling for\\Encrypted Video KQI Prediction at the Edge}
\author{Flavian Raithel \and Tobias Ho{\ss}feld \and Frank Loh\\[1mm]
\normalsize Chair of Communication Networks, University of W{\"u}rzburg, Germany}
\date{}

\begin{document}
\maketitle

\noindent
This document accompanies the paper of the same title. It collects material that does not fit
the page limit: robustness checks for the energy measurements, and the full hyperparameter
search ranges. Section and table references of the form ``Section~III-C'' point into the paper.

The measurement data and analysis code behind every number here, together with detailed
documentation of the features, the selected configurations, the measurement protocol, the
evaluation caveats and the energy model, are in the artifact repository that accompanies the
paper.

\section{Robustness Checks}\label{app:robustness}

Table~\ref{tab:retention} compares each measurement series with and without the energy-based
interquartile-range filter used in the main analysis. W1 batch-32 covers the 80 \rpi{} and 96
\jn{} configurations with valid outputs, W2 the 21/22 valid candidates and their reference
series, and W3 the four traffic views at 100 streams with and without dissection. The 96
measured configurations per device correspond to the 95 distinct hyperparameter settings of
Section~III-C, since one search selects the same trial for both selection roles.

\begin{table}[H]
\centering
\caption{Effect of energy-outlier filtering: excluded/total runs and the maximum absolute change
in a series mean relative to the unfiltered mean.}
\label{tab:retention}
\footnotesize
\setlength{\tabcolsep}{6pt}
\begin{tabular}{@{}lrr@{}}
\toprule
Workload & \rpi{} & \jn{} \\
\midrule
W1, batch 1 & 123/2880; 1.33\,\% & 58/2880; 0.88\,\% \\
W1, batch 32 & 83/2400; 6.27\,\% & 200/2880; 5.12\,\% \\
W2, paired series & 5/1260; 0.23\,\% & 19/1320; 0.17\,\% \\
W3, 100 streams & 4/240; 0.22\,\% & 2/240; 0.01\,\% \\
\bottomrule
\end{tabular}
\end{table}

For unbatched W1, removing the filter changes median whole-device savings from 71.8\,\% to
71.7\,\% on the \rpi{} and leaves the \jn{} result at 19.4\,\%. Subtracting the twelve-hour idle
baseline, $E-P_{\mathrm{idle}}T$, gives median savings of 75.8\,\% and 28.6\,\%; these are
savings per completed W1 job and do not remove the baseline of an always-on monitor. For W2,
filtering changes any candidate-minus-reference mean by at most \SI{1.87}{\joule} (\rpi{}) and
\SI{2.14}{\joule} (\jn{}).

\begin{table}[H]
\centering
\caption{Spearman correlation between the workstation energy objective of the selected trials and
their W1 whole-device edge energy (16 configurations per family and device).}
\label{tab:proxytransfer}
\footnotesize
\setlength{\tabcolsep}{10pt}
\begin{tabular}{@{}lrr@{}}
\toprule
Family & \rpi{} & \jn{} \\
\midrule
RF & 0.87 & 0.40 \\
LGBM & 0.91 & 0.86 \\
FFNN & 0.25 & 0.12 \\
CNN & 0.74 & 0.56 \\
LSTM & 0.91 & 0.91 \\
TF & 0.76 & 0.63 \\
\bottomrule
\end{tabular}
\end{table}

Table~\ref{tab:proxytransfer} shows how well the workstation energy objective used during
selection transfers to the edge. LGBM and LSTM transfer consistently, FFNN poorly, and RF depends
strongly on the execution platform, where \jn{} trees run as compiled GPU tensor operations.
Workstation-based selection is therefore a useful heuristic for some families, but it does not
guarantee edge-optimal configurations. These correlations cover only the two selected trials per
search; quantifying the heuristic would require measuring the selection regret relative to
further Pareto candidates on the edge.

\section{Hyperparameter Search Spaces}\label{app:search}

\begin{table}[H]
\centering
\caption{Architecture search ranges. $2^{a:b}$ denotes powers of two with integer exponents from
$a$ to $b$; $L$ is the input sequence length.}
\label{tab:appendixsearch}
\footnotesize
\setlength{\tabcolsep}{6pt}
\begin{tabular}{@{}l>{\raggedright\arraybackslash}p{0.78\textwidth}@{}}
\toprule
Family & Search dimensions \\
\midrule
RF & Trees 10--1000 (step 10); depth $2^{2:6}$; leaves $2^{6:12}$; minimum leaf samples 1--16; feature fraction 0.1--0.8. \\
LGBM & Trees 10--5000 (step 10); depth $2^{2:8}$; leaves $2^{4:7}$; learning rate $10^{-5}$--$10^{-1}$. \\
FFNN & 1--4 layers; width $2^{5:11}$. \\
CNN & 1--3 layers; kernels $2^{4:6}$; kernel length 2--$\min(L,10)$. \\
LSTM & 1--3 layers; hidden size $2^{4:8}$. \\
TF & 1--3 layers; heads $2^{1:3}$; model width $2^{5:8}$; feed-forward width $2^{6:10}$. \\
\bottomrule
\end{tabular}
\end{table}

Neural searches use learning rates $10^{-5}$--$10^{-1}$ (log scale), training batch sizes
$2^{7:12}$, and dropout 0--0.5 in steps of 0.1; stalling positive-class weights range from 1 to
20. Multiwindow inputs use a maximum window of 5, 10, 20, or \SI{30}{\second}, with or without
feature selection; sequence models may also use a single 0.5 or \SI{1}{\second} window. Sequence
histories contain 5--60 one-second vectors in steps of five, or 10--120 half-second vectors in
steps of ten.

The machine-readable form of every selected configuration, including its trial number, both
objective values and its feature set, is \texttt{data/features/configurations.csv} in the
artifact repository, documented in \texttt{CONFIGURATIONS.md}.

\end{document}
